\section{MVP inmutable, gossip y Day 1/100/1000}

El MVP inicial de Vercel era sumamente acotado: una llamada a la API/CLI devolvía una deployment URL inmutable. Al principio, el sistema subyacente podía estar sobredimensionado y ser costoso, pero esa interfaz bastaba para validar si los usuarios realmente querían que «desplegar equivalga a obtener una URL». Rauch tomó como fuentes de inspiración las immutable data structures, la content identity de IPFS y la gossip propagation, y después dejó que las necesidades reales impulsaran la reestructuración interna.

\subsection{Primero validar el idea-market fit, después optimizar la implementación interna}

\term{immutable deployment} significa que el contenido de un despliegue, una vez generado, no se modifica en su lugar; cada versión nueva crea una identity nueva. \term{gossip protocol} hace que los nodos se propaguen metadata entre sí hasta converger en todas las réplicas globales. El puntero del dominio puede cambiar, mientras que el deployment artifact permanece igual; así se obtienen previews trazables, promote atómico y rollback rápido.

\lecturefigure{15-immutable-mvp.png}{Una API/CLI acotada conecta la artifact identity, la propagación de metadata y la URL.}{Subtítulos locales 32:50--35:55; redibujo conceptual.}

\begin{knowledgebox}{Lectura de la figura: el MVP puede ser ineficiente por dentro, pero el contrato externo debe ser correcto}
Un sistema inicial puede hacer over-provision, siempre que el costo sea asumible y valga la pena validar la señal del producto; pero si la API externa expone al usuario el estado mutable de los servidores, migrar después resulta muy difícil. Un buen MVP elige una abstracción estable que siga siendo válida en el futuro y luego sustituye gradualmente la implementación interna.
\end{knowledgebox}

\subsection{Day 1, Day 100 y Day 1000 son tres criterios de aceptación distintos}

La subsección anterior mostró que el MVP puede ser ineficiente por dentro de forma temporal, pero que el contrato externo debe ser correcto a largo plazo; a continuación hay que usar escalas de tiempo para comprobar si ese contrato realmente absorbe la complejidad. Day 1 evalúa el onboarding y la primera impresión; Day 100 evalúa si la plataforma deja escapar complejidad cuando el proyecto crece; Day 1000 evalúa el uptime, la recuperación, el costo y la confianza de la organización. Muchas herramientas para desarrolladores solo optimizan la demo del Day 1, pero devuelven al usuario la carga del estado a largo plazo, las dependencias y la operación; por eso la madurez no puede medirse únicamente por la velocidad del primer despliegue.

\lecturefigure{16-day-maturity.png}{Un producto de infraestructura debe superar por separado las pruebas de adopción, absorción de complejidad y operación a largo plazo.}{Subtítulos locales 36:00--37:05; redibujo conceptual.}

\begin{importantbox}{Preguntas de aceptación para Day 1/100/1000}
\begin{itemize}
  \item Day 1: ¿puede una persona recién llegada desplegar de forma segura sin conocer los detalles subyacentes?
  \item Day 100: cuando aumentan las routes, los teams, los environments y las dependencias de datos, ¿sigue siendo claro el modelo predeterminado?
  \item Day 1000: ante actualizaciones, incidentes, rotación de personal y cambios de proveedor, ¿puede el sistema recuperarse y conservar su capacidad de auditoría?
\end{itemize}
\end{importantbox}

\begin{warningbox}{Un despliegue inmutable no equivale a datos inmutables}
El código y los assets pueden ser immutable, pero las bases de datos, las colas y los sistemas externos siguen cambiando. Hacer rollback de la application version puede toparse con una schema incompatibility. La plataforma necesita una migration policy, un backward-compatible contract y data recovery; no puede aplicar mecánicamente una experiencia al estilo de Git a todo el estado.
\end{warningbox}

\subsection{Resumen de la sección}

Un buen MVP fija la abstracción externa correcta y admite ineficiencias internas temporales; un buen producto de infraestructura, además, sigue superando el Day 100 y el Day 1000, absorbe la complejidad del proyecto y mantiene la confianza en su operación.
